\section{Related Work}
\label{sec:related}

\subsection{Source Identity and Privacy API Recognition}
Accurate source/sink specifications are critical for mobile privacy analysis. Android studies have mined permission mappings and source/sink lists via Stowaway, PScout, SuSi, Axplorer, and API-documentation analyses~\cite{stowaway,pscout,susi,axplorer,fuse}, and audited privacy behavior by connecting code, policies, and permissions~\cite{appaudit,appintent,policheck,ppchecker}. These works show that source/sink definitions encode platform semantics and strongly affect reported results.

\mytool applies this principle to OpenHarmony, where package migration from \texttt{@ohos.*} to consolidated \texttt{@kit.*} exports compounds the challenge. Its matching function separates API identity from exact package spelling while retaining namespace constraints. HapFlow's released \texttt{Json2ArkMethodSignature} loader resolves namespace/method entries to SDK declarations, but its rule schema does not encode executable fields or detector-side factory-origin witnesses~\cite{hapflow}. This is a limitation of the released configuration path, not of method signatures in general: a signature analysis with precise receiver types can represent manager methods. ArkAnalyzer provides \arkts code representation and call-graph construction~\cite{arkanalyzer} but does not define privacy API identity resolution. \mytool adds that evidence-carrying layer.

\subsection{Asynchronous Semantics and Lifecycle Modeling}
Mobile analysis must model lifecycle, callbacks, and implicit edges precisely. Android ICC modeling includes ComDroid, CHEX, Epicc, IC3, IccTA, DidFail, and RAICC~\cite{comdroid,chex,epicc,ic3,iccta,didfail,raicc}; EdgeMiner recovers implicit framework callbacks~\cite{edgeminer}. Hybrid-app analyses bridge Java, JavaScript, and WebView semantics~\cite{lee16,bridgescope,babelview,eoeidroid}. These studies show that missing implicit edges dominate mobile analysis errors.

OpenHarmony has analogous but different callback problems: ArkUI declarative builders, class-field callbacks, FunctionType/ClosureType values, and Promise handlers create edges absent from ordinary calls. WEMINT~\cite{wemint} demonstrates explicit callback modeling for asynchronous mini-program APIs via taint rules; \mytool differs by resolving callbacks from ArkAnalyzer's FunctionType/ClosureType IR and by handling ArkUI builder-option and receiver-factory semantics. HapFlow handles closure semantics~\cite{hapflow}; \mytool adds a bounded, separately provenance-labeled post-IFDS recovery for unresolved callback, Promise \texttt{.then()}, and \texttt{await} paths, plus report-CG edges for field and builder callbacks. The lifecycle state machine (Section~\ref{subsec:lifecycle}) captures ability/component phase transitions that the flat DummyMain cannot express.

\subsection{Evidence Representation and Large-scale Auditing}
Android taint analysis has a rich history. TaintDroid~\cite{traintdroid} motivated smartphone privacy monitoring; FlowDroid~\cite{arzt14} introduced context-/flow-/field-/object-sensitive lifecycle-aware analysis; Amandroid, DroidSafe, IccTA, LeakMiner, and AndroidLeaks explored inter-procedural data-flow, vetting, and market-side detection~\cite{wei2018amandroid,gordon2015droidsafe,iccta,leakminer,androidleaks}; HornDroid and COVERT added formal and compositional perspectives~\cite{horndroid,covert}. More recently, PacDroid unifies pointer analysis, inter-procedural propagation, ICC, and lifecycle handling and evaluates the resulting soundness--precision--performance tradeoff against several Android frameworks~\cite{pacdroid}. These systems are important methodological references, but they consume Android bytecode and Android framework models; running them on ArkTS would change both the input language and the analysis query. We therefore do not present their published numbers as an empirical baseline for \mytool.

\mytool uses the author-released HapFlow artifact as the directly comparable \openharmony taint baseline and adds recognition and subgraph evidence around it. This choice is task based rather than reputation based: HapFlow accepts ArkTS projects and answers the same source-to-sink reachability query. ArkAnalyzer and its context-sensitive Pointer Analysis Kit (APAK) provide the source IR, call graphs, and points-to information on which a taint client can be built~\cite{arkanalyzer,arktsPointer}, but do not themselves emit privacy source-to-sink findings. HarmonyFlow instead analyzes compiled Ark Panda IR and evaluates pointer/call-edge recovery~\cite{harmonyflow}; its input representation and output query differ from ours. We therefore discuss these systems as strong front-end alternatives and design influences, not as fabricated end-to-end baselines. TaintMini's Universal Data Flow Graph~\cite{taintmini} is a core analysis intermediate representation that drives cross-page data-flow resolution in mini-programs; \mytool's evidence subgraph differs in that it is a post-hoc audit artifact assembled from detection, call-chain, sink, and taint results with provenance tracking, rather than a structure that drives the analysis algorithm itself.

Malware and behavior research uses permissions, API calls, graphs, and models to classify malicious apps~\cite{drebin,droidapiminer,droidsift,apposcopy,appcontext,droidranger,droidmoss,onw19,yan12,copperdroid,intellidroid,harvester,triggerscope}; AndroZoo enables ecosystem-scale studies~\cite{Androzoo}. \mytool follows this empirical principle: the information-flow subgraph is more structured than a flat feature vector because auditors need source snippets, call contexts, sink classes, and metadata.

\mytool builds on IFDS/IDE foundations~\cite{reps95,sagiv96} and reusable infrastructures (Soot, WALA, Heros~\cite{soot,wala,heros}). ArkAnalyzer plays this role for \arkts~\cite{arkanalyzer}. APAK reports that ArkTS closure and framework modeling materially changes downstream false positives relative to CHA/RTA~\cite{arktsPointer}; this motivates our PTA ablation and our separate reporting of virtual-dispatch errors. Other \openharmony efforts explore testing~\cite{haptest,hmtest}, repair~\cite{haprepair}, and compiled-IR analysis~\cite{harmonyflow}, showing that OpenHarmony requires ecosystem-specific analysis~\cite{openharmony,harmonyroadmap,huaweiHarmonyWhitepaper}. \mytool follows the same layered principle: reuse the framework, then add domain-specific recognition and provenance-annotated subgraph mapping.

\subsection{Evaluation Methodology for Mobile Analyses}
Comparative studies show that mobile-analysis results are easily distorted by different source/sink lists, different queries, incomplete benchmark subsets, and undocumented ground truth. ReproDroid evaluates six Android taint analyzers under a common query and execution framework~\cite{reprodroid}. TaintBench goes further by storing expected and unexpected real-world flows in a machine-readable format and validating labels through multi-expert review~\cite{taintbench}. Mutation-based evaluation complements fixed benchmarks by systematically exposing unsound analysis choices~\cite{muse}. These studies shape our protocol: we use the complete author-released HapBench suite, publish a case-level oracle, keep the source and sink query fixed, count incomplete executions separately, report both positive- and negative-case metrics, and retain every misclassified case. The high-output real-project audit is reported as a precision-oriented audit rather than an independent recall benchmark because projects were selected using \mytool's output.
